\chapter{Validation Design}
\label{ch:validation}

Your validation scheme is the instrument that measures everything else. An
uncalibrated instrument does not merely add noise --- it systematically
selects for models that exploit its specific flaws. This chapter builds the
walk-forward protocol we used, explains the purged and gapped variants used
by the reference solutions, and covers the meta-question that beginners
skip: how to know whether your validation \emph{correlates with the
leaderboard at all}.

\section{Random splits are wrong here, and why}

The i.i.d.\ cross-validation reflex --- shuffle rows, split $k$ ways ---
fails on temporal panels for three separate reasons, worth keeping distinct:

\begin{enumerate}
\item \textbf{Temporal correlation}: adjacent rows are near-duplicates
  (our target is a 20-step moving average --- rows 3 steps apart share 85\%
  of their window). Random splits put near-copies of training rows into
  validation, inflating scores.
\item \textbf{Drift} (\S\ref{sec:drift}): the deployed model predicts the
  \emph{future}, where feature--target relationships have moved. Random
  splits score interpolation; the task is extrapolation.
\item \textbf{Protocol mismatch}: deployment here is a stream with daily
  label delivery and (optionally) online refitting. A static random-split
  score measures a model that will never exist in production.
\end{enumerate}

The fix is walk-forward evaluation: train on a contiguous past, validate on
a contiguous future, with an embargo at the boundary
(\S\ref{sec:embargo}).

\section{The walk-forward protocol, exactly as we run it}

Our standard evaluation (\code{run\_cv} in the repository) reproduces the
8th-place solution's protocol so that numbers are comparable across the
whole campaign:

\begin{enumerate}
\item \textbf{Split by dates, not rows.} Training pool: dates
  $[d_0, d_1]$; validation tail: dates $(d_1, d_2]$, typically 100 days.
  A one-day embargo separates them.
\item \textbf{Fit} the model (and all preprocessing) on the pool.
\item \textbf{Evaluate day by day, in order.} For each validation day $d$:
  first, if online refitting is enabled, take one small gradient step on
  day $d{-}1$'s now-revealed labels (this mirrors the competition's lag
  delivery exactly); then predict all of day $d$ one timestep at a time
  (in practice: causally, so that no within-day future is visible).
\item \textbf{Score} the concatenated predictions with the exact
  competition metric --- weights included, zero-mean form included.
\end{enumerate}

Two scores come out of every run: \textbf{static} (step 3 without the refit)
and \textbf{online}. Report both, always. The gap between them is itself a
diagnostic: large for recurrent models ($+0.003$ to $+0.005$ here), near
zero for attention models, catastrophically negative when the refit learning
rate is mis-set (Chapter~\ref{ch:online}) --- three different stories that a
single blended number would hide.

\begin{keybox}[: validate the process, not the checkpoint]
If deployment includes adaptation, validation must include adaptation. A
model is not a frozen function; it is a \emph{procedure} --- weights, plus
update rule, plus preprocessing state. Our biggest single scoring lever
(online refit) is invisible to any validation scheme that scores frozen
checkpoints.
\end{keybox}

\section{Purged and gapped $k$-fold: when you need more than one fold}

A single train/tail split uses data inefficiently and yields one noisy
number. The standard upgrade for temporal data is \emph{purged group
$k$-fold with embargo} --- the scheme the DRW winner used (six contiguous
two-month groups, gap of one group):

\begin{itemize}
\item Partition time into $k$ contiguous groups.
\item Each fold validates on one group, trains on the others, and
  \textbf{purges} training rows whose label windows overlap the validation
  group, plus an embargo margin on each side.
\item Aggregate scores across folds --- and \emph{report the spread}, not
  just the mean. The DRW winner's per-fold Pearson correlations were
  $\{0.196, 0.160, 0.034, 0.142, 0.146, 0.111\}$: one regime (fold 3) was
  nearly unpredictable. That spread is information --- it tells you score
  variance across regimes, which a single tail cannot.
\end{itemize}

Note what purged $k$-fold does \emph{not} give you: folds that train on the
future and validate on the past are still extrapolating backward, which is
optimistic under drift. For final decisions we trust forward-only splits;
purged $k$-fold earns its keep for \emph{feature selection} (fold-consistency
screening, Chapter~\ref{ch:selection}) where you want several verdicts
per candidate.

\section{The gap variant for honest deployment simulation}

The competition's private test period began months after the training data
ended. A validation tail adjacent to training therefore \emph{overstates}
deployable performance: your model enjoys a freshness the real one will not
have. The fix is a \textbf{gap}: train on $[d_0, d_1]$, skip $g$ days,
validate on $[d_1{+}g, d_2]$. The 8th-place solution ran both gapless and
200-day-gap variants; the gap costs score (staleness) and the \emph{size of
that cost} estimates your drift exposure on the private set. Models that
degrade least across the gap --- typically the online-refitted ones, which
re-adapt during the validation walk --- are exactly the ones to trust for a
delayed test period.

\section{Calibrating the instrument: CV--LB correlation}

The most important property of a validation scheme is not its absolute
level but whether \emph{improvements transfer}. The DRW winner stated his
CV was far from the public leaderboard in level but moved with it
consistently --- ``an increase in my local CV leads to a consistent increase
in public LB'' --- and that correlation, not the level, is what he relied
on. Protocol for establishing it:

\begin{enumerate}
\item Submit 3--5 models spanning a range of local scores (deliberately
  include a weak one; you need dynamic range).
\item Plot local score against leaderboard score. You are looking for
  monotonicity, not identity.
\item If monotonic: trust local, spend submissions sparingly, ignore
  level differences. If not: your local protocol mismatches the test
  distribution --- check the information set, the metric implementation,
  the time period, and per-day weighting before trusting anything again.
\end{enumerate}

Budget submissions like the scarce measurements they are: each one is a
draw from the leaderboard's own noise distribution, and public LBs of
low-SNR competitions are themselves noisy enough that chasing them
\emph{is} overfitting (the final-standing shakeups of every such
competition are the receipts).

\section{Validation-set bookkeeping at low SNR}

Three quantitative habits, restated from Chapters~\ref{ch:problem}
and~\ref{ch:leakage} because they live here operationally:

\begin{itemize}
\item \textbf{Know your measurement noise.} Re-score the same model with
  different seeds ($\pm 0.0003$--$0.0008$ for our nets) and across two
  validation windows. Effects below this floor are unresolved, full stop.
\item \textbf{Two-tier discipline}: an iteration window and a
  finalists-only window. The second window's verdicts are spent, not
  reused.
\item \textbf{Log everything} --- every run's config, score, and seed into
  an append-only artifact directory. Selection bias thrives on forgotten
  losers; the log is how you remember that today's ``breakthrough'' idea
  already lost twice under different names (our repository's
  \code{artifacts/bench/} plus \code{WHAT\_DIDNT\_WORK.md} serve exactly
  this function).
\end{itemize}

\section{Choosing the validation period}

Under drift, \emph{when} you validate is a modeling decision:

\begin{itemize}
\item \textbf{Recent tail} (our default: the last 100 days) --- closest to
  the test distribution; the honest choice for final decisions.
\item \textbf{Multiple windows} --- the stability gate of
  \S\ref{sec:atlas}: any structural claim must reproduce across distant
  windows. Feature findings, refit learning rates, and window-length
  choices all went through this gate.
\item \textbf{Regime coverage} --- if the training span includes an obvious
  regime break (ours does: day-length stabilizes at date 677, and early
  dates behave differently enough that the reference solution simply
  discards dates before 700), validate within the regime you will deploy
  into, and consider training there too
  (Chapter~\ref{ch:preprocessing}).
\end{itemize}

\begin{jsbox}[: the protocol as a fixed contract]
Every number in this book's tables was produced by the same contract:
train through date 1598, one-day embargo, walk the tail 1599--1698 day by
day with online refit where stated, score with the exact weighted zero-mean
$R^2$. Freezing the contract early --- resisting the temptation to
``improve'' it mid-campaign --- is what makes a hundred experiments
mutually comparable, and comparability is what turns experiments into
knowledge.
\end{jsbox}

\begin{drillbox}
Take a model you have already trained on any temporal dataset and score it
three ways: random 80/20 split, forward split, forward split with a gap
equal to 10\% of your data span. The spread between the three numbers is
your dataset's own measurement of temporal correlation plus drift --- and a
preview of how much a naive validation scheme would have lied to you.
\end{drillbox}
